%-------------------------------------------------------------------------------
\section{Normative Rule Extraction}
\label{sec:rule-extraction}
%-------------------------------------------------------------------------------

This section presents DLEP, our framework for extracting Web PKI normative rules into a structured intermediate representation.

\subsection{PKI Knowledge Graph Builder}
\label{sec:method-kg-builder}

The PKI knowledge graph builder is the offline component that converts heterogeneous Web PKI normative sources into the retrieval substrate used by DLEP.
It normalizes documents into a hierarchical structure with stable identifiers for specifications, sections, headings, and text units.
It then extracts explicit references such as section citations, RFC references, and policy cross-links, and links them to target nodes when the targets can be resolved unambiguously.
In parallel, it records certificate-field concepts and aliases, namely subjectAltName, dNSName, basicConstraints.cA, and policy OIDs, to map natural-language mentions to canonical certificate paths during later extraction.

PKI specifications form a densely interconnected cross-document network.
For example, RFC~5280 references multiple RFCs, while CABF~BR directly depends on RFC~5280 field definitions.
We model this knowledge as a multi-edge directed graph with eight node types, namely Specification, Section, Definition, CertificateField, Rule, Operation, Value, and Concept, and four source-backed relations, namely CONTAINS, DEFINES, REFERENCES, and APPLIES\_TO. Appendix~\ref{app:kg} summarizes these relations and their deterministic retrieval constraints.

This builder is deliberately conservative.
It stores only relations directly supported by document structure or textual evidence, while leaving normative judgments such as conflict, override, and lintability to later deterministic stages.
Normative-judgment relations such as CONFLICTS\_WITH and OVERRIDES are produced independently by the rule engine and are not included in the retrieval graph.
The resulting graph serves as a traceable context index rather than as an oracle. It helps the extractor identify definitions, inherited constraints, and referenced rules, but does not determine whether a requirement is lintable.

\subsection{Deterministic Context Retrieval}
\label{sec:method-graphrag}

To address cross-document references, DLEP organizes context deterministically rather than relying on the LLM to infer relations.
We implement a deterministic retrieval module based on GraphRAG~\cite{edge2024local} that treats the knowledge graph of \S\ref{sec:method-kg-builder} as a fixed input and retrieves context through bounded subgraph queries.

Given a target Section node, the retrieval module performs a BFS expansion over the $k$-hop neighborhood using the source-backed CONTAINS, DEFINES, REFERENCES, and APPLIES\_TO relations.
It then assembles LLM context by relation-type priority and provides term definitions, field metadata, and referenced sections. The retrieval process uses neither vector similarity nor LLM inference, and produces no independent judgment.

\subsection{DLEP}
\label{sec:method-dual-layer}

\begin{table}[tb]
\centering
\caption{RFC~2119 keyword tiers.}
\label{tab:rfc2119}
\setlength{\tabcolsep}{4pt}
\begin{tabularx}{\columnwidth}{@{}L{0.27\linewidth} L{0.30\linewidth} X@{}}
\hline
Tier & Keywords & Semantics \\
\hline
Mandatory      & MUST, SHALL, REQUIRED   & Absolute requirement; a violation is non-compliance \\
Prohibition    & MUST NOT, SHALL NOT     & Absolute prohibition \\
Recommendation & SHOULD, SHOULD NOT, RECOMMENDED & Strongly recommended or discouraged; justified deviation permitted \\
Optional       & MAY, OPTIONAL           & Fully optional \\
\hline
\end{tabularx}
\end{table}

\begin{table}[tb]
\centering
\caption{IR core four-tuple fields.}
\label{tab:ir-core}
\setlength{\tabcolsep}{4pt}
\footnotesize
\begin{tabularx}{\columnwidth}{@{}L{0.18\linewidth} X@{}}
\hline
Field & Meaning \\
\hline
subject    & Certificate field path, for example extensions.\allowbreak subjectAltName.\allowbreak dNSName, resolved by FieldResolver \\
obligation & RFC~2119 keyword, such as MUST/SHALL/SHOULD/\dots \\
predicate  & Assertion type, such as must\_be\_present, conform\_to, equal, \dots \\
constraint & constraint value or pattern \\
\hline
\end{tabularx}
\end{table}

To address the LLM controllability challenge, we design DLEP as a dual-layer extraction pipeline that preserves the LLM's semantic understanding while keeping the extraction process auditable.

Layer~1 performs deterministic recall.
Web PKI specifications commonly express normative requirements using the RFC~2119 keywords listed in Table~\ref{tab:rfc2119}.
RFC~8174 further specifies that these keywords carry normative force only when capitalized, enabling deterministic keyword matching.

To maximize recall, Layer~1 operates in three passes.
Pass~1 directly matches RFC~2119 keywords.
Pass~2 detects nested structures so subordinate rules inherit the parent obligation level.
Pass~3 captures normative statements that do not follow the standard RFC~2119 form.

Layer~2 performs controlled semantic parsing using an LLM, but restricts the model to a constrained language-understanding role.
Instead of generating lint rules or compliance judgments directly, Layer~2 converts normative text into a JSON-serialized intermediate representation, abbreviated IR, while all normative reasoning is deferred to later deterministic stages.

The IR decouples semantic extraction from downstream tasks, allowing the same extraction result to support multiple uses, as shown in Fig.~\ref{fig:ir-downstream}.
Each IR record is serialized as a JSON object with 33 fields.
Its core component, shown in Table~\ref{tab:ir-core}, is the four-tuple
$\mathrm{IR} = \langle\text{subject}, \text{obligation}, \text{predicate}, \text{constraint}\rangle$,
together with 29 extension fields summarized in Appendix~\ref{app:ir}.

This IR-based design has three advantages over direct end-to-end generation of lint checks from raw text.
First, citation, constraint, and obligation information are stored separately, improving traceability and conflict handling.
Second, the IR serves as deterministic input to the four-condition lintability analysis in \S\ref{sec:method-lintability}, making lintability decisions reproducible rather than model-dependent.
Third, once a rule is determined to be lintable, generating static lint checks from the IR is more stable and interpretable because key elements such as field paths, assertion types, and constraint values are already explicitly structured.

To constrain LLM behavior, Layer~2 enforces a predefined schema and category set.
Every extracted rule must map to explicit source text, and certificate field hierarchies must follow the ASN.1 path definitions in RFC~5280.
The IR generation process is also staged. The model first determines the rule category and then fills in the remaining fields.
Appendix~\ref{app:ir} provides the category definitions, correction rules, and tuning details.

\begin{figure}[tb]
  \centering
  \paperfig{figures/fig02_ir_downstream.png}{IR-driven ``extract once, apply
    many'' architecture.}
  \caption{IR-driven ``extract once, apply many'' architecture.}
  \label{fig:ir-downstream}
\end{figure}
